\chapter{Literature Review}
\label{chap:litreview}

\Cref{chap:intro} established that a mature body of fairness research has not yet converged into a reusable auditing platform. This chapter surveys that body of research in the specific areas that shape BiasAperture's design: the formal definitions behind its Core~Four metrics, the critiques that justify pairing every disparity with a significance test, the finding that motivates its explainability layer, the conventions that shape its report structure, the benchmark datasets against which it is validated, and the regulatory frameworks its output is designed to evidence.

\section{Foundational Bias Discovery}

The empirical case for facial-analysis auditing rests on Buolamwini and Gebru's Gender Shades study, which showed that commercial gender-classification services performed markedly worse on darker-skinned female subjects than on lighter-skinned male subjects \cite{buolamwini2018gendershades}, discussed already in \cref{chap:intro}. Dehdashtian et al.'s subsequent survey situates that single finding within a wider taxonomy, cataloguing where bias enters the computer-vision pipeline, data collection, annotation, training, and deployment, and the corresponding families of fairness metrics and mitigation techniques the field has proposed in response \cite{dehdashtian2024survey}. Their survey is the clearest evidence that the metrics and mitigation literature is mature while the auditing-workflow literature is not: a rich taxonomy of \emph{what} to measure exists, but comparatively little standardises \emph{how} to measure it repeatably.

\section{Formal Fairness Metrics}

Two of BiasAperture's Core~Four metrics, \gls{eod} and \gls{eop}, are direct operationalisations of definitions introduced by Hardt et al. \cite{hardt2016equality}. Equalized odds is the stricter of the two: it requires parity in both the true-positive rate and the false-positive rate across demographic groups, so that the model performs equally well for every group regardless of the true label. Equal opportunity relaxes this to parity in the true-positive rate alone, for the positive outcome class specifically, prioritising equitable allocation of favourable outcomes over symmetric error rates across both classes. BiasAperture computes both rather than choosing one, since the two definitions answer different fairness questions and neither is a general substitute for the other; \gls{dpd} and \gls{dir} then supply the outcome-rate and ratio-based perspectives Hardt et al.'s error-rate framing does not itself cover.

\section{Critiques of Naive Fairness Thresholds}

A disparate-impact ratio is only as trustworthy as the threshold used to interpret it, and Watkins et al. show this is a genuine hazard rather than a theoretical one \cite{watkins2022fourfifths}. Their critique of the U.S. four-fifths rule argues that importing an 80\% threshold from employment-discrimination law into algorithmic fairness without its original statistical and sociotechnical context is an instance of what they term epistemic trespassing, and that treating the resulting ratio as a binary pass or fail can grant a false sense of legal and ethical safety while conflating compliance with actual fairness. This finding is the direct justification for FR-004 and NFR-001: BiasAperture never reports \gls{dir} as a bare ratio against an implicit threshold, but always alongside a chi-squared test of independence and a bootstrap confidence interval, so that a disparity is characterised by its statistical strength rather than measured against a borrowed and contested cut-off.

\section{Proxy Variables and Explainability}

Kurian et al.'s study of bias encoding in convolutional networks trained on medical images provides the finding that motivates \cref{chap:methodology}'s explainability layer: a \gls{cnn} can learn to exploit demographic information even when no demographic label is present anywhere in its training data, by encoding it as a proxy through unrelated, seemingly neutral visual features such as texture or lighting \cite{kurian2024bias}. Applied to facial analysis, this means a detected subgroup disparity cannot be assumed to arise only from the demographic label the audit is stratifying on; it may instead be an artefact of a correlated visual feature the model has entangled with that label. This is precisely why FR-005 requires \gls{shap}-based attribution on every flagged disparity rather than treating a Core~Four result as self-explanatory: the explainability layer exists to help distinguish a genuinely demographic effect from an unrelated confound before the finding is reported as evidence.

\section{Reporting and Documentation Conventions}

BiasAperture's report structure is not invented from scratch. Mitchell et al.'s Model Cards propose a standardised, structured summary of a model's intended use, performance characteristics, and evaluation results, aimed explicitly at making that information legible to both technical and non-technical stakeholders \cite{mitchell2019modelcards}. Gebru et al.'s Datasheets for Datasets propose the complementary discipline for the data side: documenting a dataset's motivation, composition, collection process, and recommended uses so that downstream users can reason about its limitations \cite{gebru2018datasheets}. FR-006's requirement that every generated report be structurally informed by both conventions means BiasAperture's dataset-provenance and model-performance sections follow documentation patterns the field already recognises, rather than an ad hoc format each reader must learn from scratch.

\section{Benchmark Datasets}

BiasAperture is validated principally against FairFace, a face-attribute dataset explicitly constructed for balanced representation across race, gender, and age, using a seven-category race taxonomy designed to reduce the classification skew common in earlier face datasets \cite{karkkainen2021fairface}. UTKFace is used as a secondary benchmark, but its age annotations are model-estimated via a regression approach rather than human-labelled, a known source of label noise that \cref{chap:intro}'s Limitations section already flags and that the platform's risk register treats explicitly: every subgroup metric is reported with its sample size and confidence interval precisely so that a result resting on noisier labels is not presented with unwarranted certainty.

\section{The Regulatory-Technical Bridge}

The EU AI Act (Regulation 2024/1689) and the NIST AI RMF, both introduced in \cref{chap:intro}, state \emph{what} must be governed and measured, respectively, but neither specifies how a development team turns that obligation into a concrete technical check. Buscemi et al. close that gap directly, decomposing the AI Act's high-risk-system obligations into a stable taxonomy of high-level requirements and sub-requirements and mapping each to a concrete, technology-agnostic technical verification activity \cite{buscemi2025assessing}, an operational-verification approach that reduces the interpretive uncertainty a development team otherwise faces when translating a legal article into a specific test. \Cref{sec:regmapping} follows the same operational spirit at a scale appropriate to a single Article~10 mapping: each Core~Four metric, each significance test, and each sample-size guard is tied to the specific data-governance clause it evidences, rather than the report gesturing at compliance in general terms.

\section{Synthesis}

No single artefact in the literature surveyed above combines all of the following: a reusable software pipeline, rather than a one-off manual audit; the Core~Four metrics computed by cross-validating, independently maintained toolkits rather than a single implementation; every disparity paired with the significance testing Watkins et al.'s critique shows is necessary rather than a bare threshold; \gls{shap}-based diagnosis of the proxy-variable risk Kurian et al. identify; a report structured against the Model Cards and Datasheets conventions; and explicit, clause-level traceability to the EU AI Act and NIST AI RMF in the operational spirit Buscemi et al. establish. \Cref{chap:requirements,chap:methodology} specify exactly this combination as BiasAperture's design, closing the engineering gap identified in \cref{chap:intro}.
